\documentclass[12pt]{article}

\usepackage{sbc-template}

\usepackage{graphicx,url}

\usepackage[brazil]{babel}
\usepackage[utf8]{inputenc}

\usepackage{float}
\usepackage{tikz}
\usetikzlibrary{arrows.meta,positioning}

\sloppy

\title{WhyRunner: Online Judge for Competitive Programming}

\author{Juan Israel dos Anjos}

\address{Instituto Federal de Educacao, Ciencia e Tecnologia de Minas Gerais
  (IFMG)\\
  Rua Afonso Sardinha, 90, Bairro Minas Talco – Ouro Branco/MG – CEP: 36494-018 – Brazil
  \email{juaniwk3@gmail.com}
}

\begin{document}

\maketitle

\begin{abstract}
  Existing online judges primarily focus on correctness-based evaluation, often neglecting pedagogical feedback and real-time contest dynamics essential for educational contexts. This paper introduces WhyRunner: Programming Online Judge, a modular and secure system that integrates containerized code execution, real-time contest management, a professor-curated lesson/exercise authoring model, and AI-powered feedback. This paper presents the system's architecture, development methodology, and implementation; classroom evaluation is left to future work.
\end{abstract}

\begin{resumo}
  Os juízes online existentes focam principalmente na avaliação baseada em correção, negligenciando o feedback pedagógico e a dinâmica de competições em tempo real, essenciais em contextos educacionais. Este artigo apresenta o WhyRunner: Programming Online Judge, um sistema modular e seguro que integra execução de código em contêineres, gerenciamento de competições em tempo real, um modelo de autoria de aulas/exercícios curado pelo docente e feedback assistido por inteligência artificial. Este artigo descreve a arquitetura do sistema, a metodologia de desenvolvimento e a implementação; a avaliação em sala de aula fica como trabalho futuro.
\end{resumo}

\section{Introduction}

Competitive programming plays a significant role in algorithmic thinking, logical reasoning, and problem-solving skills. These capabilities are essential for students pursuing careers in computer science and related fields. In recent years, the adoption of online judge systems (automated platforms that evaluate submitted code against predefined test cases) has become increasingly popular, especially in academic settings. These platforms simulate real-world programming contests, enabling students to practice and improve their skills outside of traditional classroom environments.

Most online judges are not optimized for educational use. They were primarily designed for competitive environments, prioritizing fairness and efficiency over pedagogical considerations. As a result, students receive limited feedback typically restricted to ``Wrong Answer'' or ``Time Limit Exceeded'' without explanations that help learning.

Many platforms are rigid in design, making it difficult for educators and students to configure and use. Features such as adaptive problem sets, real-time contest control, and support for multiple programming languages are either limited or absent. The increasing complexity of software development tools and the diverse backgrounds of students further highlight the need for more accessible and educationally driven platforms.

To address these limitations, this work presents \textbf{WhyRunner}, a modern and modular online judge system developed specifically for educational contexts. WhyRunner was designed from the ground up with three core objectives: (1) to provide rich, automated feedback using artificial intelligence; (2) to support safe, scalable, and multi-language code execution through containerization; and (3) to enable flexible real-time contest management for classroom or remote settings.

This paper details the design, development, and current implementation status of WhyRunner.

\section{Methodology}\label{sec:methodology}

The system development followed an iterative and incremental methodology. This means the system was built step by step, adding features and improvements after each test and evaluation.

\subsection{Requirements Analysis}

The initial phase consisted of identifying shortcomings in existing judges through a combination of literature review, informal conversations with students and instructors, and the author's own practical experience with competitive programming platforms. Key problems found included limited pedagogical feedback, lack of support for different programming languages and website display languages, and contest management.

\subsection{Design Principles}

The system was designed with modularity, scalability, and pedagogical value as central principles:

\begin{itemize}
  \item \textbf{Modularity}: Components such as code execution, feedback generation, user interface, and contest management were developed as loosely coupled modules, enabling easier maintenance and future upgrades.
  \item \textbf{Scalability}: The system leverages message queues and container orchestration to handle multiple concurrent submissions with minimal latency.
  \item \textbf{Security}: All submitted code is executed in isolated Docker containers with restricted resource access, minimizing the risk of malicious code affecting the system.
  \item \textbf{Educational Feedback}: Instead of reporting only binary results (e.g., ``Accepted'' or ``Wrong Answer''), the system uses AI to generate contextual explanations and suggestions for improvement.
\end{itemize}

\subsection{Technology Evaluation}

The selection of technologies for WhyRunner was guided by practicality, familiarity, and empirical testing. Throughout development, multiple frameworks and tools were evaluated based on performance, maintainability, and ease of deployment in cloud environments.

Initially, the front-end was built using the React framework due to prior experience and its robust ecosystem. To ensure server-side rendering, route handling, and production-readiness, the Next.js framework was selected. It offered a seamless development experience, especially for projects requiring both front-end and server-side logic.

Queue management went through two discarded approaches before settling: BullMQ (Redis-based) generated excessive idle polling traffic in production, and its replacement, AWS SQS, added external operational overhead. The final architecture uses PostgreSQL's native \texttt{LISTEN/NOTIFY} instead, eliminating the external queue entirely while keeping asynchronous, event-driven behavior.

Similarly, Supabase was dropped in favor of a standalone PostgreSQL instance (via Drizzle ORM) once it proved restrictive for custom server-side logic, and a separate NestJS API was dropped in favor of Next.js Server Actions once the split-repository model made shared types and auth logic too costly to maintain. A short-lived TurboRepo monorepo was also abandoned, as it complicated Docker-based deployment of the judge worker, which now ships as a standalone Rust binary listening for submissions over \texttt{LISTEN/NOTIFY} and executing them in Docker.

\subsection{Development Process}

Development occurred in cycles, each involving implementation, integration, informal testing, and refinement. A minimum viable product (MVP) was shown in a Discord server to a small group of students for informal feedback outside any classroom setting. That feedback led to rapid improvements, including enhancements to the user interface, better error reporting, and more detailed AI explanations. Each subsequent cycle refactored code, adjusted parameters (e.g., time/memory limits), and improved usability based on the author's own testing.

\subsection{Current Validation Status}

WhyRunner has not yet been deployed or evaluated in a real classroom setting. All development-cycle feedback so far came from informal testing (the Discord MVP session above) and the author's own use of the system; no student or instructor cohort, contest, or usability questionnaire has been run against it. A classroom evaluation is planned as future work (\S\ref{sec:future-work}).

\section{Related Work} \label{sec:related}

\subsection{Ataux} \cite{silva2021ataux}

Ataux lets an administrator create programming problem lists that a class submits solutions to, generating a results sheet. It offers admin/collaborator/member roles, difficulty/type filtering, team registration, and solving via Codeforces, on a RESTful React/NestJS/PostgreSQL stack deployed on Heroku. Its supported-language range is limited compared to larger platforms, and feedback is correctness-only. Usability was assessed with PSSUQ across 71 students and 1 teaching assistant over two weeks.

\subsection{GOJ} \cite{lima2022gamma}

GOJ stores and automatically judges problems from past UnB Programming Marathon editions, organized by event, with beginner tutorials and admin tooling (Gamma Judge Admin). It is a microservices/Web-API architecture on AWS (SQS, SNS, S3), .NET, and MongoDB, with a React front-end and Bash-script-based judging.

\subsection{Run Codes} \cite{runcodes2024}

Run Codes is an online judge platform developed by Brazilian universities, designed to create and manage programming contests and activities. It allows teachers to create problem lists and competitions, automates the evaluation of code submissions, and generates reports for performance analysis.

It supports a wide variety of programming languages and integrates with virtual learning environments. However, Run Codes' complexity and external dependencies can make it less flexible for fully customized learning activities or research projects.

\subsection{BOCA} \cite{de2004boca}

BOCA (BOCA Online Contest Administrator) is an open-source system widely used in Brazilian programming competitions like the ACM ICPC regional contests. BOCA allows managing contests, teams, problems, and automated judging.

The system is built mainly with PHP and relies on a LAMP stack (Linux, Apache, MySQL, PHP). While robust for traditional contests, its interface and extensibility are more limited compared to newer platforms.

\section{Comparative Analysis}

Comparing the analyzed platforms to WhyRunner:

\begin{table}[h]
\centering
\begin{tabular}{|l|c|c|c|c|}
\hline
\textbf{System} & \textbf{AI Feedback} & \textbf{Multi-Language} & \textbf{Real-time Features} & \textbf{Curated Lessons} \\
\hline
Ataux      & No  & Limited   & No   & No \\
GOJ        & No  & Moderate  & Limited & No \\
Run Codes  & No  & Extensive & Yes  & No \\
BOCA       & No  & Moderate  & Yes  & No \\
WhyRunner  & Yes & Extensive & Yes  & Yes \\
\hline
\end{tabular}
\caption{Comparison of programming judge systems}
\label{tab:judge-comparison}
\end{table}

Beyond AI feedback, WhyRunner's second differentiator is the professor-curated Class/Lesson/Exercise model (\S\ref{sec:class-lesson-model}): none of the compared systems let an instructor assemble an ordered, gradable assignment from existing problems and review it per exercise outside of a timed contest. Solution-constraint enforcement (\S\ref{sec:prepublish-validation}) is a related but separate differentiator: it lets a professor require a specific algorithmic approach, not just a correct output, which none of the compared systems support either.

\section{System Design}

\subsection{Overview}

WhyRunner is structured as two cooperating services: a \textbf{Next.js web application} responsible for the user interface, authentication, and all server-side business logic via Server Actions; and a \textbf{Rust-based judge worker} responsible for code execution and result evaluation.

The system's backend relies on \textbf{PostgreSQL} for persistent data storage, managed by \textbf{Drizzle ORM}. Asynchronous communication between the web application and the judge worker is handled through PostgreSQL's native \texttt{LISTEN/NOTIFY} mechanism, eliminating the need for an external message broker. When a user submits code, the web application inserts a row into the \texttt{submission} table and issues a \texttt{NOTIFY}; the judge worker, which holds an open connection via \texttt{PgListener}, receives the notification and begins processing. The \textbf{Google Gemini API} generates AI-powered feedback for students on failed submissions.

Beyond timed contests, WhyRunner also supports a professor-curated Class/Lesson/Exercise model: a professor assembles lessons from existing problems for their class, and students submit and are graded exercise-by-exercise, independent of the contest scoreboard flow (\S\ref{sec:class-lesson-model}).

\subsection{Technology Stack}

\begin{itemize}
    \item \textbf{Next.js} \cite{nextjs2024}: Full-stack web framework for the user interface and server-side logic via Server Actions.
    \item \textbf{Better-Auth:} Authentication library integrated with Next.js for session management.
    \item \textbf{Drizzle ORM} \cite{drizzleorm2024}: Type-safe database access layer for PostgreSQL.
    \item \textbf{PostgreSQL} \cite{stonebraker1986design}: Relational database for persistent data storage and inter-service messaging via \texttt{LISTEN/NOTIFY}.
    \item \textbf{Rust} \cite{matsakis2014rust}: Systems language used for the judge worker, chosen for performance and memory safety.
    \item \textbf{Docker} \cite{merkel2014docker}: Isolated, resource-constrained containers for safe user code execution.
    \item \textbf{Gemini API} \cite{gemini2023}: Powers AI-driven feedback on failed submissions.
\end{itemize}

\newpage
\subsection{Submission Pipeline}\label{sec:submission-pipeline}

The core implementation revolves around a pipeline that ensures decoupling between submission, execution, constraint checking, and feedback. Figure~\ref{fig:submission-flow} illustrates the current flow, which now branches by submission type (contest/standalone vs. lesson exercise) and includes the solution-constraint checks introduced in \S\ref{sec:prepublish-validation}.

\begin{figure}[H]
\centering
\begin{tikzpicture}[
  node distance=0.5cm and 1.6cm,
  box/.style={draw, rounded corners, fill=gray!10, text width=4.6cm, align=center, minimum height=0.9cm, font=\footnotesize},
  note/.style={draw, dashed, rounded corners, fill=gray!5, text width=4.0cm, align=center, minimum height=0.8cm, font=\scriptsize},
  arr/.style={-{Latex[length=2.2mm]}, thick}
]

\node[box] (submit) {Student submits code (contest, standalone, or lesson exercise)};
\node[box, below=of submit] (save) {Server Action validates, saves submission as \texttt{PENDING}, issues \texttt{NOTIFY}};
\node[box, below=of save] (pickup) {Judge worker (\texttt{PgListener}) picks up the submission and its test cases};
\node[box, below=of pickup] (exec) {Docker container compiles and runs the code against each test case};
\node[note, right=of exec] (policy) {Contest/standalone: stop at first failing case. Lesson exercise: run every case, for a per-exercise pass fraction};
\node[box, below=of exec] (constraints) {If the exercise has solution constraints, run structural static analysis, then (if configured) an LLM classification for the algorithm requirement};
\node[note, right=of constraints] (statuses) {Outcomes: \texttt{PASSED}, \texttt{FAILED}, \texttt{ERROR}, or \texttt{CONSTRAINT\_VIOLATION}};
\node[box, below=of constraints] (persist) {Final status written to PostgreSQL; in-app notification created};
\node[box, below=of persist] (client) {Frontend updates via TanStack Query; student sees results, may request an AI hint on failure};

\draw[arr] (submit) -- (save);
\draw[arr] (save) -- (pickup);
\draw[arr] (pickup) -- (exec);
\draw[arr] (exec) -- (constraints);
\draw[arr] (constraints) -- (persist);
\draw[arr] (persist) -- (client);
\draw[arr, dashed] (exec) -- (policy);
\draw[arr, dashed] (constraints) -- (statuses);

\end{tikzpicture}
\caption{Current submission pipeline: grading policy branches by submission type, and solution-constraint checks run before the final status is persisted}
\label{fig:submission-flow}
\end{figure}

\subsubsection{Sending Submission}
The user submits code via the web interface. A Next.js Server Action validates the submission, checks contest membership and rate limits, and saves it to PostgreSQL with a \texttt{PENDING} status. A \texttt{NOTIFY} is then issued on the \texttt{submission\_created} channel.

\subsubsection{Queue Dispatch}
The judge worker holds a persistent connection to PostgreSQL and listens on the \texttt{submission\_created} channel using \texttt{PgListener} from the \texttt{sqlx} library. Upon receiving a notification, it fetches the submission and its associated test cases from the database.

\subsubsection{Execution and Grading Policy}
The judge worker spawns a Docker container for the submission's language, with strict resource limits (\texttt{--cpus 0.5 --memory 128m}) and a configurable execution timeout. Contest and standalone submissions stop at the first failing test case; lesson-exercise submissions run every declared test case, so the pass count reflects exactly how many cases the code got right (\S\ref{sec:class-lesson-model}).

\subsubsection{Constraint Checking}
If the exercise has active solution constraints, the checks described in \S\ref{sec:prepublish-validation} run once I/O grading completes: structural static analysis first, then, if configured, an algorithm-requirement classification via Gemini. Either violation path resolves the submission to \texttt{CONSTRAINT\_VIOLATION} instead of \texttt{PASSED}.

\subsubsection{AI Feedback}
On failed submissions, the student may request AI assistance. The system calls the Gemini API with a structured prompt containing the problem statement, the student's code, and the judge's failure report. The AI is instructed to provide vague but directional hints, supporting learning rather than providing direct solutions.

\subsubsection{Result Notification}
The final status and optional feedback are persisted. The front-end polls for updates via TanStack Query, and students can view their submission history with detailed results per test case. Beyond polling, the system raises an in-app notification whenever a submission finishes grading, alongside notifications for social (follow, like, comment), contest (join request/approval), and lesson (newly unlocked, assignment-related) events; each user has a notification inbox with an unread count and read/unread state, and repeated same-type events (e.g. several likes on one post) are aggregated into a single running-count entry rather than flooding the inbox.

\subsection{Class, Lesson, and Exercise Model}\label{sec:class-lesson-model}

Alongside timed contests, WhyRunner offers a second, asynchronous mode built on three entities: a \textbf{Class}, a \textbf{Lesson} (a professor-curated assignment scoped to one class), and an \textbf{Exercise} (an existing problem attached to a lesson, in a fixed order). A professor creates a class, assembles a lesson from existing problems, sets an optional due date, and publishes it; students join via a code and submit code per exercise through the same judge pipeline used elsewhere. A per-lesson toggle can reveal each exercise's expected test-case outputs (inputs are always visible), trading answer-secrecy for easier manual review. Students submit the whole lesson at once; after the due date, the professor reviews it and leaves free-text feedback per exercise. Submitting with exercises still unsolved requires confirmation, so partial submissions are always intentional.

Grading differs from the contest path in one respect: the judge runs every declared test case for a lesson-exercise submission instead of stopping at the first failure, so each exercise's score is derived automatically as the fraction of test cases passed, and a lesson's total score sums those fractions, weighting every exercise equally regardless of its test-case count.

\subsection{Execution Environment}

To ensure safe and reproducible code execution, user code runs inside Docker containers spawned by the Rust judge worker. Each container is ephemeral: it is created per submission, executed with time and resource limits, and destroyed upon completion. Resource constraints are enforced via Docker runtime flags:

\begin{itemize}
  \item \textbf{CPU}: limited to 0.5 cores (\texttt{--cpus 0.5})
  \item \textbf{Memory}: limited to 128 MB (\texttt{--memory 128m})
  \item \textbf{Execution time}: enforced by the judge worker via a configurable timeout
\end{itemize}

The system currently supports Python, Rust, and C++ as execution languages. Each language environment includes an interpreter or compiler and a test runner script. Support for additional languages can be added by extending the runner configuration.

\subsection{AI-Assisted Authoring and Feedback}

The Gemini API backs four distinct capabilities in WhyRunner, one student-facing and three professor-facing, all following the same pattern: a structured prompt with user-supplied content enclosed in XML delimiters to prevent prompt injection, and a result that is either returned directly to the requester or left as an editable draft that must be explicitly saved before it affects the problem.

\subsubsection{Student Hint Feedback}

On a failed submission, a student may explicitly request AI help rather than receiving it automatically, supporting learner autonomy and limiting API usage. The \texttt{getAIHelp} Server Action sends the problem statement, the student's code and language, and the judge's failure report (status, failing input, expected vs. actual output) to Gemini, with an instruction to stay vague about the direct fix while explaining the logical error; user-supplied content is XML-delimited against prompt injection, and the response is localized to the user's language. This reduces teacher workload while giving contextual, constructivist-aligned assistance instead of a binary verdict.

\subsubsection{AI-Assisted Problem Authoring}

Three further, non-mutating capabilities target the problem creator from the edit workspace, each returning an editable draft rather than writing directly:

\begin{itemize}
  \item \textbf{AI-drafted reference solution:} from a title/topic, proposes a description, example I/O, and a reference solution for pre-publish validation (\S\ref{sec:prepublish-validation}).
  \item \textbf{AI problem review:} critiques the draft's description and suggests edge-case test cases (input, output, rationale), each addable with one click.
  \item \textbf{AI-generated narrative:} short, optional themed flavor text framing the problem, never required for publishing.
\end{itemize}

\subsection{Pre-Publish Validation and Solution Constraints}\label{sec:prepublish-validation}

Two mechanisms let a professor set stricter grading guarantees than ordinary test-case matching.

\textbf{Pre-publish I/O validation.} A draft problem's reference solution runs through the same sandboxed judge execution used for student submissions, against the problem's current test cases, before the problem can be published; validation runs are stored separately and never appear in leaderboards or history. A passing run is tied to the exact test-case content (and any constraints, below) it ran against, so editing either after a pass marks it stale and re-blocks publish.

\textbf{Solution constraints.} Attached per exercise, constraints come in two kinds: \emph{structural} rules from a fixed catalog (max loop-nesting depth, forbidden/required constructs), checked via lightweight source-text scanning across all six submission languages rather than a real parser, which can false-positive on a token inside a string or comment; and a single optional \emph{algorithm-requirement} constraint, a free-text description (e.g. ``must use Dijkstra's algorithm'') checked once test cases and structural rules pass, via a Gemini 2.5 Flash call given only the requirement text, submission language, and source (never the problem statement, since correctness was already verified), returning a structured \texttt{\{satisfied, rationale\}} result. Either kind of violation resolves the submission to a distinct \texttt{CONSTRAINT\_VIOLATION} status, separate from \texttt{PASSED}/\texttt{FAILED}/\texttt{ERROR}, denying contest and completion credit even though every test case passed, with the violated rule (and, for algorithm requirements, the model's rationale) shown to the student.

\subsection{Integration and Deployment}

The system is organized as a two-project repository: the Next.js web application and the Rust judge worker. Both services share the same PostgreSQL instance and communicate exclusively through database primitives, with no shared runtime dependencies.

The web application is deployed on a platform supporting long-running Next.js applications. The judge worker is built as a Docker image and deployed as a persistent background service, maintaining its \texttt{LISTEN} connection to PostgreSQL at all times.

\section{Evaluation Status}

WhyRunner has not yet been evaluated with a real cohort of students or instructors: no classroom contest, lesson assignment, or usability study has been run against the system as described in this paper. What follows is a description of the system as implemented and manually exercised by the author during development, not a report of a controlled trial. A classroom evaluation, and the quantitative usability study it would enable, is planned future work (\S\ref{sec:future-work}).

\section{Discussion}

\subsection{Strengths}

WhyRunner exhibits several notable strengths:

\begin{itemize}
  \item \textbf{Modular and scalable architecture:} Clear separation of concerns across the web layer, judge worker, and execution containers allows for easy maintenance and future expansion.
  \item \textbf{Secure code execution:} User code runs in ephemeral Docker containers with enforced CPU, memory, and time limits, isolating it from the host system.
  \item \textbf{No external queue dependency:} Using PostgreSQL \texttt{LISTEN/NOTIFY} for inter-service communication removes operational overhead while maintaining reliability.
  \item \textbf{AI-powered feedback:} The use of large language models enhances the learning process by offering tailored hints and performance improvement suggestions without revealing direct solutions.
  \item \textbf{Type-safe full-stack:} The combination of Next.js Server Actions, Drizzle ORM, and Zod validation provides end-to-end type safety from database schema to user interface.
\end{itemize}

\subsection{Limitations}

Despite its strengths, WhyRunner has some limitations:

\begin{itemize}
  \item \textbf{Dependency on third-party AI APIs:} The feedback system relies on the availability and consistency of the Gemini API, which may be subject to rate limits or cost constraints at scale.
  \item \textbf{Limited language support:} Currently, Python, Rust, and C++ are fully supported for execution. C and Java are architecturally supported but not yet implemented in the judge runner.
  \item \textbf{Deployment complexity:} The initial setup requires knowledge of Docker, Rust, and PostgreSQL configuration, which may present a barrier for non-technical instructors.
  \item \textbf{No classroom validation yet:} As noted in \S\ref{sec:methodology}, WhyRunner as a whole, including contests, the Class/Lesson/Exercise model (\S\ref{sec:class-lesson-model}), AI-assisted authoring, and solution-constraint enforcement (\S\ref{sec:prepublish-validation}), is functionally complete but has not been evaluated with a real cohort of students or instructors.
  \item \textbf{Structural constraint analysis is heuristic, not a real parser:} Loop-nesting and forbidden/required-construct checks (\S\ref{sec:prepublish-validation}) scan raw source text (indentation width or brace depth, plus word-boundary substring search) rather than building an AST, so a matching token inside a string literal or comment can produce a false positive.
\end{itemize}

\subsection{Future Work}\label{sec:future-work}

Planned improvements include:

\begin{itemize}
  \item \textbf{Classroom evaluation:} Deploy WhyRunner, both the contest path and the Class/Lesson/Exercise model, in a real classroom setting with students and an instructor, and apply standardized usability questionnaires (e.g., SUS or PSSUQ) to collect measurable data on satisfaction and usability across the full feature set, including AI-assisted authoring and solution-constraint enforcement.
  \item \textbf{Language support expansion:} Add support for C, Java, and Kotlin by implementing the respective compiler scripts in the Docker execution environment.
  \item \textbf{Plagiarism detection:} Integrate code similarity analysis to support academic integrity in contest and lesson settings.
  \item \textbf{Improved test case management:} Develop tooling for automated test case generation and validation based on problem constraints, building on the AI-assisted edge-case suggestions already available in problem review.
\end{itemize}

\section{Conclusion}

WhyRunner introduces a novel approach to online code evaluation in educational settings by combining real-time contest features, secure multi-language execution, and AI-powered feedback. It has grown beyond timed contests into a professor-curated Class/Lesson/Exercise authoring platform, adding AI-assisted problem drafting, review, and narrative generation; pre-publish reference-solution validation; solution-constraint enforcement; per-exercise grading; and an in-app notification system, extending the same secure judge pipeline to asynchronous, instructor-led coursework. The system as described here is implemented and self-tested by the author, but not yet evaluated with real students or instructors; that classroom evaluation is the immediate next step (\S\ref{sec:future-work}).

With a modular and extensible design, WhyRunner can evolve to meet broader needs in computer science teaching and competitive programming training. The continued integration of AI technologies and developer tooling positions it as a promising platform for the future of automated learning environments.

\bibliographystyle{sbc}
\bibliography{refs}

\end{document}
